\section{Animasi (animation) dan Keyframes}

\textbf{CSS Animation} memungkinkan urutan perubahan properti yang didefinisikan melalui \code{@keyframes}. Berbeda dari transition yang memerlukan pemicu (hover/focus), animation dapat berjalan otomatis saat halaman dimuat, berulang tanpa batas, dan memiliki kontrol arah (maju, mundur, bergantian). Animation cocok untuk loading spinner, attention-seekers, dan efek masuk (entrance effects) \cite{mdnCss}.

\begin{webcode}[language=CSS, caption={Keyframes dan animation}]
@keyframes slideIn {
  from { transform: translateX(-100%); opacity: 0; }
  to   { transform: translateX(0);     opacity: 1; }
}
@keyframes pulse {
  0%, 100% { transform: scale(1); }
  50%      { transform: scale(1.05); }
}
.card   { animation: slideIn 0.5s ease forwards; }
.badge  { animation: pulse 2s ease infinite; }
\end{webcode}

\begin{table}[htbp]
\centering
\begin{tabular}{|P{5cm}|P{3.5cm}|P{3.5cm}|}
\hline
\textbf{Properti} & \textbf{Fungsi} & \textbf{Contoh} \\
\hline
\code{animation-name} & Nama keyframe & \code{slideIn} \\
\hline
\code{animation-duration} & Durasi satu siklus & \code{0.5s} \\
\hline
\code{animation-iteration-count} & Jumlah pengulangan & \code{1}, \code{infinite} \\
\hline
\code{animation-direction} & Arah pemutaran & \code{normal}, \code{alternate} \\
\hline
\code{animation-fill-mode} & State setelah selesai & \code{forwards}, \code{backwards} \\
\hline
\end{tabular}
\caption{Properti animation CSS yang sering digunakan}
\end{table}

\begin{catatan}
\textbf{Transition vs Animation:}
\begin{itemize}
  \item \textbf{Transition}: perlu pemicu (hover/focus/class change); satu arah; cocok untuk interaksi pengguna.
  \item \textbf{Animation}: berjalan otomatis; dapat berulang dan multi-tahap (keyframes); cocok untuk efek dekoratif.
  \item Untuk respek aksesibilitas, periksa \code{prefers-reduced-motion} dan kurangi/matikan animasi bagi pengguna yang memilih pengurangan gerak \cite{mdnAccessibility}.
\end{itemize}
\end{catatan}

